\section{Filesystem State Space and Invariants}\label{sec:state}

We define the object on which the recovery operator acts. The
mathematical scaffolding of this section is preserved from the v1
report~\cite{desbrieres2026ftrfs}: the state space, the layout
abstraction, the consistency relation, and the five invariants
\textbf{I1}--\textbf{I5} are unchanged. What changes between v1
and v2 is the perturbation family quantified over by the
soundness theorem in \cref{sec:soundness}; the layout
substrate on which that theorem is stated is the same.

\subsection{State Space}

\begin{definition}[Layered state]\label{def:state}
The BEAMFS state space~\(\mathcal{S}\) is the Cartesian product
\[
  \mathcal{S} \;=\; \mathcal{S}_V \times \mathcal{S}_P,
\]
where~\(\mathcal{S}_V\) is the volatile component (kernel data
structures, page cache, in-memory superblock copy, in-memory
bitmap) and~\(\mathcal{S}_P\) is the persistent component (on-disk
superblock, inode table, allocation bitmap, data blocks, RS event
journal).
\end{definition}

\begin{definition}[Layout]\label{def:layout}
A \emph{layout}~\(\Lambda\) is a finite map from named regions
(\texttt{sb}, \texttt{inode\_table}, \texttt{bitmap},
\texttt{rs\_journal}, \texttt{data}) to byte ranges in their
hosting layer, together with a redundancy specification (CRC
algorithm, RS parameters) per region.
\end{definition}

The FTRFS layout~\cite{fuchs2015ftrfs,desbrieres2026ftrfs} is one
instance of~\(\Lambda\). BEAMFS v2 instantiates an extended
layout: in the INLINE scheme (s\_data\_protection\_scheme =
\textsc{universal\_inline}), every region is RS-protected, and
each data block is partitioned into sixteen RS subblocks
\(\mathrm{RS}(255, 239)\) carrying user content plus parity
symbols. The on-disk format is normatively specified in the
companion document \texttt{Documentation/format-v4.md} of the
implementation~\cite{beamfs-impl}.

\subsection{Consistency Relation}

\begin{definition}[Per-region consistency]
Let~\(R\) be a region of~\(\Lambda\) with redundancy
specification~\(r\). A state~\(s\) is \emph{R-consistent}, written
\(s \models_R\), iff applying~\(r\) to the bytes of~\(R\) in~\(s\)
yields a zero syndrome (CRC valid and Reed--Solomon syndromes all
zero across all subblocks of \(R\)).
\end{definition}

\begin{definition}[Consistency]\label{def:consistent}
A state~\(s = (s_V, s_P) \in \mathcal{S}\) is \emph{consistent},
written~\(s \models\), iff:
\begin{enumerate}
\item For every region~\(R\) of~\(\Lambda_V\) (the volatile-side
  layout), \(s \models_R\).
\item For every region~\(R\) of~\(\Lambda_P\) (the persistent-side
  layout), \(s \models_R\).
\item For every region~\(R\) that has both volatile and persistent
  representations (the superblock, the bitmap, recently accessed
  inodes), the canonical serialization of~\(s_V|_R\) equals
  \(s_P|_R\) byte-for-byte (\emph{coherence}).
\item Every structural sentinel (the zero pads of \texttt{re\_reserved}
  and \texttt{re\_pad} in the v4 RS event journal entry, the
  validity range of \texttt{s\_data\_protection\_scheme}, the high
  zero bytes of feature flags) is at its canonical value.
\end{enumerate}
We write~\(\bot\) for the explicit fail-closed state, distinct
from every consistent state.
\end{definition}

\subsection{The Five Invariants}\label{sub:invariants}

The recovery soundness theorem (\cref{thm:soundness-v21}) requires
the following invariants to hold of the layout~\(\Lambda\) and of
the operator~\(\mathcal{G}\). They are minimal and stated
separately so that an implementation can be audited against them
line by line.

\begin{description}
\item[I1 -- Total redundancy.] Every region of~\(\Lambda\) is
  covered by a redundancy specification (CRC for detection, RS for
  correction) whose correction power equals or exceeds the maximum
  per-subblock symbol-error budget targeted by the layout. For
  \(\mathrm{RS}(n, k)\) with \(n - k\) parity symbols and
  per-subblock correction capacity \(t = \lfloor (n-k)/2 \rfloor\),
  the layout's stated EMR perturbation family must contain only
  events whose per-subblock symbol-error count does not exceed
  \(t\). For BEAMFS v2 INLINE: \(\mathrm{RS}(255, 239)\) with
  \(t = 8\) per subblock and 16 subblocks per 4\,KiB block, the
  per-block budget is 128 symbol errors.
\item[I2 -- Sentinel separation.] Each region carries at least one
  byte whose canonical value is fixed and whose corruption is
  detectable in the redundancy syndrome alone, without semantic
  interpretation.
\item[I3 -- Coherence checkpoint.] For any region present in both
  layers, the canonical serialization map is total, deterministic,
  and bijective on its domain of valid values. This rules out
  ambiguous in-memory representations of an on-disk byte sequence.
\item[I4 -- Idempotence.] For every consistent~\(s\),
  \(\mathcal{G}(s) = s\). The recovery operator is the identity on
  consistent states.
\item[I5 -- Fail-closed monotonicity.] If
  \(\mathcal{G}(s) = \bot\) for some~\(s\), then
  \(\mathcal{G}(\bot) = \bot\). Once the operator declares
  fail-closed, no subsequent application reverses that decision.
\end{description}

The first three are properties of the layout; the last two are
properties of the operator. The soundness theorem couples them
under the EMR family.

\begin{remark}
The space overhead of \(\mathrm{RS}(n, k)\) is \((n-k)/k\). For
\(\mathrm{RS}(255, 239)\), this is 6.7\%. The INLINE scheme
applies this overhead per 4\,KiB disk block (sixteen
\(\mathrm{RS}(255,239)\) subblocks of 255 bytes each cover the
4\,096-byte block plus parity), with 4\,080 bytes of user content
per block. The space overhead is preserved across the v1 and v2
designs; what differs is per-subblock granularity (16 subblocks
per block versus FTRFS's per-region granularity), enabling
recovery up to \(t \cdot 16 = 128\) symbol errors per block.
\end{remark}

\begin{remark}
The v1 report acknowledged that I1 is unrealistically strong as a
universal requirement; the present report retains that
acknowledgement. A future v3 may explore selective redundancy
under a probabilistic relaxation of the soundness statement (with
bounded silent-corruption risk quantified rather than zero by
construction). The conservative v2 statement is what allows the
soundness theorems of \cref{sec:soundness} to be stated
\emph{absolutely} on the EMR-bounded family, not statistically.
\end{remark}
